% Claim proof file for row claim_3_12 -- "HardInterventionNodeSplitOrder".
% Auto-generated stub by scaffold/solving_current_row.py at row start. A
% manager agent dedicated to writing the TeX proof fills in the body below.
% The proof must:
%   - Stay close to the lecture notes' proof if one exists (always search
%     the LN first for a `\begin{proof}` block immediately following the
%     claim; copy / lightly edit when present).
%   - Otherwise use the LN paradigm and cite prior claims by their `ref`.
%   - Be self-contained enough that a mathematician can follow it without
%     re-reading the LN.
%   - Have no `sorry`, `omitted`, or "obvious" shortcuts.
%
% Compiles standalone via the `subfiles` package.

\documentclass[main]{subfiles}

\begin{document}
% ref: claim_3_12 (proof, with statement restated for readability)
% title: HardInterventionNodeSplitOrder
\def\rowref{claim\_3\_12}
\phantomsection\label{claim_3_12}
%
% Cross-references: see claim_<ref>_statement_<title>.tex in this folder
% for the corresponding statement.

% --- statement (restated) ---------------------------------------------------
\begin{Rem}[HardInterventionNodeSplitOrder]
    Note that if $W_1$ and $W_2$ are not disjoint and $w \in W_1 \cap W_2 \ins V$
    then first hard intervening on $w$ turns $w$ into an input node, for now indicated as $w^i$, and a node-splitting hard intervention (if we would define it for input nodes) would not change $w^i$.
    If, on the other hand, we would first split the node $w$ into $w^o$ and $w^i$ then we would first need to resolve the ambiguity on which of those two the hard intervention should be applied. A hard intervention on $w^i$ would not do anything, but would leave the additional output node $w^o$ in the graph, while hard intervening on $w^o$ would turn $w^o$ into an additional input node, for now indicated as $(w^o)^i$. So in the latter case we are left with two input node $(w^o)^i$, which does not have any edges, and $w^i$, which might have outgoing edges.
\end{Rem}

% --- proof ------------------------------------------------------------------
\begin{proof}
\medskip\noindent\emph{Preamble.}
The Remark walks through what happens when the disjointness premise of
\refrow{claim_3_11} fails -- that is, when there is some
$w \in W_1 \cap W_2 \ins V$ shared between the hard-intervention target
$W_1$ and the SWIG target $W_2$. We unpack its content as three operational
scenarios plus a punchline corollary, mirroring the LN's three-clause
narrative: \emph{(A)} ``HI first, then SWIG'' fails to define a CDMG at
all; \emph{(B1)} ``SWIG first, then HI on the input copy $w^i$'' is a
no-op; \emph{(B2)} ``SWIG first, then HI on the output copy $w^o$''
genuinely shrinks the output set; and the \emph{punchline} that, in the
singleton case $W_1 \cap W_2 = \lC w \rC$, branches (B1) and (B2) deliver
non-equal CDMGs. Throughout we let $G = (J, V, E, L)$ be a CDMG and follow
the SWIG conventions of \refrow{def_3_12}: for the SWIG argument $W \ins V$,
each $w \in W$ has a fresh output copy $w^o$ and a fresh input copy $w^i$
with $w^o \neq w^i$, while for $v \in (J \cup V) \sm W$ we use the
identification $v^o = v^i = v$. As in the proof of \refrow{claim_3_11}
each CDMG is presented as a 4-tuple $(J_\ast, V_\ast, E_\ast, L_\ast)$.

\medskip\noindent\emph{Scenario A: ``HI first, then SWIG'' is ill-typed.}
Fix $w \in W_1 \cap W_2$ with $w \in V$. Consider the left-hand side of
\refrow{claim_3_11}'s commute identity,
$\lp G_{\doit(W_1)} \rp_{\swig(W_2)}$. By \refrow{def_3_12}, the SWIG
operator $\swig(W_2)$ is defined on a host CDMG only when its target $W_2$
is a subset of the host's output set, so this composite requires
$W_2 \ins V_{\doit(W_1)}$. But \refrow{def_3_10}.ii gives
$V_{\doit(W_1)} = V \sm W_1$, and the assumption $w \in W_1 \cap W_2$
yields $w \in W_2$ together with $w \in W_1$, so $w \notin V \sm W_1$.
Therefore $W_2 \not\ins V \sm W_1 = V_{\doit(W_1)}$ and the SWIG
precondition fails: $\lp G_{\doit(W_1)} \rp_{\swig(W_2)}$ cannot be formed
as a CDMG in this framework.

This matches the LN's prose verbatim: ``first hard intervening on $w$ turns
$w$ into an input node, for now indicated as $w^i$'' is the unfolding
$J_{\doit(W_1)} = J \cup W_1 \ni w$, $V_{\doit(W_1)} = V \sm W_1 \not\ni w$
from \refrow{def_3_10}.i--ii; after this promotion $w$ is no longer
available as an output to feed into $\swig(W_2)$. The LN's parenthetical
``a node-splitting hard intervention (if we would define it for input
nodes) would not change $w^i$'' is an aside acknowledging that $\swig$ in
the LN has the same domain restriction $W \ins V$ that ours does, so even
a hypothetical extension to input nodes would only redundantly preserve
$w^i$. The formal content of Scenario A is therefore precisely the
typing obstruction recorded above. \checkmark

\medskip\noindent\emph{Scenario B1: SWIG first, then HI on input copies is a no-op.}
Take any $W \ins V$ and any subset $W' \ins W$, and let
$(W')^i := \lC w^i \st w \in W' \rC \ins W^i$ denote the corresponding set
of input copies inside the SWIG $G_{\swig(W)}$. We claim
\[
  \lp G_{\swig(W)} \rp_{\doit((W')^i)} \;=\; G_{\swig(W)}.
\]
Write $G' := G_{\swig(W)}$ for brevity, and recall from
\refrow{def_3_12} that
\[
  J' = J \dcup W^i, \qquad V' = (V \sm W) \dcup W^o,
\]
with $W^i \cap V' = \emptyset$ by the fresh-copies clause of
\refrow{def_3_12} (the LN's ``we consider $w^o \neq w^i$ for $w \in W$''
combined with $V \sm W$ being disjoint from the fresh $W^i$). Note also
that every directed edge of $G'$ has the form $v_1^i \tuh v_2^o$ for some
$v_1 \tuh v_2 \in E$ (so its \emph{head} $v_2^o$ lies in $V'$, since
$v_2 \in V$ implies $v_2^o \in W^o$ if $v_2 \in W$ and $v_2^o = v_2 \in V \sm W$
otherwise), and every bidirected edge of $G'$ has the form
$v_1^o \huh v_2^o$ for some $v_1 \huh v_2 \in L$ (so both endpoints lie in
$V'$ by the same case analysis applied to $v_1, v_2 \in V$ -- using
\refrow{def_3_1}'s constraint $L \ins V \times V$). We verify the four
data fields of $\lp G' \rp_{\doit((W')^i)}$ against those of $G'$ in turn.

\emph{Input set.} By \refrow{def_3_10}.i,
$J_{(G')_{\doit((W')^i)}} = J' \cup (W')^i$. Since
$(W')^i \ins W^i \ins J'$, the union is absorbed: $J' \cup (W')^i = J'$.
\checkmark

\emph{Output set.} By \refrow{def_3_10}.ii,
$V_{(G')_{\doit((W')^i)}} = V' \sm (W')^i$. We have
$(W')^i \ins W^i$ and $V' \cap W^i = \emptyset$ (recorded above), so
$V' \cap (W')^i = \emptyset$ and $V' \sm (W')^i = V'$. \checkmark

\emph{Directed edges.} By \refrow{def_3_10}.iii,
$E_{(G')_{\doit((W')^i)}} = E' \sm \lC v \tuh u \st u \in (W')^i \rC$.
Every directed edge of $G'$ has head $v_2^o \in V'$, and
$V' \cap (W')^i = \emptyset$, so no head ever lies in $(W')^i$ and the
removed set is empty. \checkmark

\emph{Bidirected edges.} By \refrow{def_3_10}.iv,
$L_{(G')_{\doit((W')^i)}} = L' \sm \lC v \huh u \st u \in (W')^i \rC$.
Every bidirected edge of $G'$ has both endpoints in $V'$, so neither
endpoint lies in $(W')^i \ins W^i$ (disjoint from $V'$). Again the
removed set is empty. \checkmark

All four components of $\lp G' \rp_{\doit((W')^i)}$ coincide with those
of $G' = G_{\swig(W)}$, so the two CDMGs are equal. This is the formal
content of the LN's ``A hard intervention on $w^i$ would not do
anything'': inside the SWIG, every input copy in $W^i$ is already an
input vertex (in $J'$, not in $V'$) and is not incident to any edge,
because directed edges target $V'$ and bidirected edges have both
endpoints in $V'$. Promoting elements of $W^i$ ``from input to input''
via $\doit$ therefore makes no change at all. In particular, taking
$W' = \lC w \rC$ specialises this to the singleton scenario of the
Remark.

\medskip\noindent\emph{Scenario B2: SWIG first, then HI on output copies shrinks $V$.}
Take any $W \ins V$ and any subset $S \ins W$, and let
$S^o := \lC s^o \st s \in S \rC \ins W^o$. We compute the output set
of $\lp G_{\swig(W)} \rp_{\doit(S^o)}$ and find that it strictly
shrinks (relative to $V_{\swig(W)}$) by exactly $S^o$.

By \refrow{def_3_10}.ii,
$V_{\lp G_{\swig(W)} \rp_{\doit(S^o)}} = V_{\swig(W)} \sm S^o$. Substituting
$V_{\swig(W)} = (V \sm W) \dcup W^o$ from \refrow{def_3_12}.ii, and noting
that $S^o \ins W^o$ is disjoint from $V \sm W$ (since $W^o$ consists of
fresh copies not equal to any $v \in V \sm W$), the difference splits as
\[
  V_{\swig(W)} \sm S^o
  \;=\; (V \sm W) \,\dcup\, \lp W^o \sm S^o \rp
  \;=\; (V \sm W) \,\dcup\, (W \sm S)^o,
\]
where the last equality uses injectivity of the map $w \mapsto w^o$ on
$W$ (the fresh-copies clause of \refrow{def_3_12}: distinct $w$ give
distinct $w^o$). Equivalently, identifying $v^o = v$ for $v \in V \sm W$
(so that $V \sm W = (V \sm W)^o$), this rewrites as
\[
  V_{\lp G_{\swig(W)} \rp_{\doit(S^o)}}
  \;=\; (V \sm S)^o,
\]
where the right-hand side denotes $\lC v^o \st v \in V \sm S \rC$. The
output set has therefore shrunk by exactly the fresh-output-copy block
$S^o$, which is non-trivial whenever $S \neq \emptyset$. \checkmark

This is the formal content of the LN's ``hard intervening on $w^o$
would turn $w^o$ into an additional input node, for now indicated as
$(w^o)^i$'': by \refrow{def_3_10}.i--ii applied to the SWIG, each $s^o$
($s \in S$) is removed from $V_{\swig(W)}$ and added to the input set --
the LN's $(s^o)^i$ -- and the directed/bidirected edges incident to
$s^o$ are deleted (by \refrow{def_3_10}.iii--iv). For the punchline we
only need the change in $V$, which is the cleanest witness; the rest of
the four-component update is not used below.

\medskip\noindent\emph{Punchline: in the singleton scenario, (B1) and (B2) yield different CDMGs.}
Fix $w \in V$ and set $W := \lC w \rC$, so that $W \ins V$ and the SWIG
$G_{\swig(\lC w \rC)}$ is well-defined. The LN's punchline contrasts two
natural ``post-SWIG hard interventions on $w$'': one targets the input
copy $\lC w^i \rC \ins W^i$, the other targets the output copy
$\lC w^o \rC \ins W^o$. We claim
\[
  \lp G_{\swig(\lC w \rC)} \rp_{\doit(\lC w^o \rC)}
  \;\neq\;
  \lp G_{\swig(\lC w \rC)} \rp_{\doit(\lC w^i \rC)}
\]
by exhibiting a vertex that lies in the output set of one side but not
the other.

Applying Scenario~B1 with $W' = \lC w \rC$ (so $(W')^i = \lC w^i \rC$)
gives
\[
  \lp G_{\swig(\lC w \rC)} \rp_{\doit(\lC w^i \rC)} \;=\; G_{\swig(\lC w \rC)},
\]
whose output set is
\[
  V_{\swig(\lC w \rC)} \;=\; (V \sm \lC w \rC) \,\dcup\, \lC w^o \rC
\]
by \refrow{def_3_12}.ii. In particular this set \emph{contains} $w^o$.

Applying Scenario~B2 with $S = \lC w \rC$ (so $S^o = \lC w^o \rC$) gives
\[
  V_{\lp G_{\swig(\lC w \rC)} \rp_{\doit(\lC w^o \rC)}}
  \;=\; (V \sm \lC w \rC)^o
  \;=\; V \sm \lC w \rC,
\]
where the second equality uses the identification $v^o = v$ for
$v \notin W = \lC w \rC$. This set does \emph{not} contain $w^o$: the
fresh output copy $w^o$ is distinct from every $v \in V \sm \lC w \rC$
(since for such $v$ we have $v^o = v \neq w^o$, the latter being a fresh
copy of $w$ that is not equal to any non-$W$ vertex).

The two output sets therefore disagree on the single vertex $w^o$,
which suffices for the two CDMGs to disagree as 4-tuples. This is the
Remark's punchline: in branch (B2) the resulting graph carries two input
copies of $w$ -- the LN's ``$(w^o)^i$, which does not have any edges''
(because $w^o$ inherited the incident edges of $w$ in $G$, all of which
$\doit(\lC w^o \rC)$ deletes by \refrow{def_3_10}.iii--iv) together with
the original SWIG input ``$w^i$, which might have outgoing edges''
(since edges $v_1^i \tuh v_2^o$ of the SWIG with $v_1 = w$ survive into
$\lp G_{\swig(\lC w \rC)} \rp_{\doit(\lC w^o \rC)}$ unchanged) -- while
branch (B1) leaves only the original $w^i$ and keeps $w^o$ in the
output set. The disagreement on $V$ is the formal record of this
asymmetry. \checkmark

\medskip
Combined, Scenarios~A, B1, B2 and the punchline show that when
$W_1 \cap W_2 \ins V$ is non-empty no order/choice of operations can
recover the commute identity of \refrow{claim_3_11}: the ``HI first''
order is ill-typed, and the ``SWIG first'' order yields genuinely
distinct CDMGs depending on whether the subsequent hard intervention
targets the input copy $w^i$ or the output copy $w^o$ of the shared
vertex $w$. The disjointness premise of \refrow{claim_3_11} is therefore
essential, not a convenience.
\end{proof}

\end{document}
